% !TeX root = ../presentacion.tex
% ===========================================================================
%  2. Cómo lo construimos — 2 min
%
%  Este bloque no lo pide la rúbrica de forma explícita: cae dentro de
%  «Calidad técnica y documentación» (10 %). Va aquí, y no al final, porque
%  enlaza directamente con la sesión 08 del curso (Vibe Coding vs.
%  Spec-Driven Development) y explica por qué el resto de la charla puede
%  enseñar trazabilidad de un requisito a una línea de código.
%
%  Si el tiempo aprieta, este es el primer bloque que se recorta a una sola
%  lámina. No el de arquitectura.
% ===========================================================================

\bloque{Cómo lo construimos}{1,5 min}{\exponeArquitectura}%
       {Calidad técnica y documentación (10\,\%)}

\begin{frame}{Especificación primero, no «vibe coding»}
  \guion{Enlazar explícitamente con la sesión del curso. El mensaje: elegimos
  el enfoque que el propio curso señala como el indicado para sistemas con
  reglas de negocio y varios equipos tocando el mismo contrato.}

  \begin{columns}[T]
    \begin{column}{0.48\textwidth}
      \begin{block}{Vibe coding}
        \footnotesize
        Se describe el destino, no la ruta. Fricción mínima, iteración
        conversacional.\\[3pt]
        Bien para: prototipos, exploración de una API, demos.
      \end{block}
    \end{column}
    \begin{column}{0.48\textwidth}
      \begin{block}{Spec-Driven Development}
        \footnotesize
        La especificación es la fuente de verdad; el código es su
        consecuencia.\\[3pt]
        Necesario cuando hay: reglas de negocio críticas, contratos entre
        equipos, trazabilidad exigible.
      \end{block}
    \end{column}
  \end{columns}

  \vfill
  \begin{center}
    Este proyecto tiene las tres cosas. \alert{Elegimos la segunda.}
  \end{center}
\end{frame}

\begin{frame}{El flujo, y lo que dejó como evidencia}
  \guion{Enseñar el camino de un requisito hasta su prueba. Es el argumento
  que sostiene la tabla de reglas de negocio del bloque 7: cada RN apunta a
  una clase y a un test con nombre propio.}

  \begin{center}
    \small
    constitución $\rightarrow$ \textbf{spec} $\rightarrow$ plan $\rightarrow$
    tasks $\rightarrow$ implementación $\rightarrow$ pruebas
  \end{center}

  \vspace{4pt}
  \begin{columns}[T]
    \begin{column}{0.5\textwidth}
      \begin{itemize}\footnotesize
        \item Contrato OpenAPI acordado \emph{antes} de escribir el servicio.
        \item Criterios de aceptación como pruebas, no como prosa.
        \item Cada regla de negocio tiene una clase y un test con nombre.
        \item Las compuertas de revisión quedan en el historial del repositorio.
      \end{itemize}
    \end{column}
    \begin{column}{0.46\textwidth}
      \begin{block}{Lo que compró}
        \footnotesize
        Poder responder «¿de dónde sale esto?» señalando un archivo, en vez
        de reconstruirlo de memoria.
      \end{block}
    \end{column}
  \end{columns}

  \note{Si el docente pregunta por herramientas: Spec Kit sobre el
  repositorio. Tenerlo listo, pero no abrir ese melón sin que lo pidan.}
\end{frame}
